\chapter{Introduction}

This blueprint accompanies the Lean~4 formalization in the
\href{https://github.com/Langlandsproject/Langlands}{Langlands} repository.
Its scope follows the design document at
\texttt{docs/superpowers/specs/2026-05-23-affine-algebraic-group-design.md},
which records twelve phases of work from \emph{affine algebraic groups} up
through the Chevalley classification of split reductive groups.

Each blueprint node carries a back-reference to the GitHub issue that tracks
its implementation. The convention for these references is illustrated below.

\begin{convention}\label{conv:typeclass-idiom}\issue{1}\leanok
Throughout this project, an \emph{affine algebraic group} $G$ over a base scheme $S$
is represented in Lean by the data
\[
  (G : \Sch) \;,\; [G.\mathsf{Over}\,S] \;,\; [\GrpObj(\asOver\,G\,S)]
\]
together with property typeclasses such as $[\mathsf{IsAffineHom}\,(G \downarrow S)]$
and $[\mathsf{LocallyOfFiniteType}\,(G \downarrow S)]$. We deliberately avoid
introducing a bundled type \texttt{GroupScheme S}; the Mathlib idiom is the
canonical reference. Two aggregator typeclasses
(Definitions~\ref{def:affine-group-scheme} and~\ref{def:algebraic-group}) are
provided as Level~1 sugar over the underlying stack.

This convention is the subject of Phase A1; the namespace and documentation
hub live in
\texttt{lean/LanglandsLean/AlgebraicGeometry/Conventions.lean}.
\end{convention}
